\chapter{Local and Weighted Estimation}
\label{cap:local}
% ── index ───────────────────────────────────────────────────────
\index{weighted least squares|textbf}
\index{WLS|see{weighted least squares}}
\index{wls@\texttt{wls()}|textbf}
\index{rolling OLS|textbf}
\index{rols@\texttt{rols()}|textbf}
\index{LOWESS|textbf}
\index{lowess@\texttt{lowess()}|textbf}
\index{local estimation|textbf}
\index{nonparametric smoothing}

% ─────────────────────────────────────────────────────────────────────────────
\section{Weights, windows and smoothing}

Classical estimation (OLS, MLE) treats each observation symmetrically.
Three alternatives allow the model to adapt locally:

\begin{description}
  \item[WLS] Weighted least squares: observations with lower variance
    (or greater relevance) receive higher weight.

  \item[Rolling OLS] Estimates the model in sliding time windows,
    revealing how coefficients evolve over time.

  \item[LOWESS] Nonparametric smoothing: for each point $x_0$, fits
    a locally weighted regression using nearby neighbors.
\end{description}

\subsection*{WLS (Weighted Least Squares)}

The WLS estimator minimizes:
\[
  \hat{\beta}_{WLS} = \arg\min_\beta \sum_{i=1}^n w_i (y_i - x_i'\beta)^2
  = (X'WX)^{-1} X'Wy
\]
where $W = \text{diag}(w_1, \ldots, w_n)$. If the weights are proportional to
$1/\sigma_i^2$ (error variances), WLS is the efficient estimator under
heteroscedasticity (= GLS for independent errors).

\subsection*{Rolling OLS}

For time series non-stationary in parameters (time-varying coefficients),
rolling OLS with window $h$ estimates:
\[
  \hat{\beta}_t = (X_{t-h:t}' X_{t-h:t})^{-1} X_{t-h:t}' Y_{t-h:t}
\]

\subsection*{LOWESS}

LOWESS (Cleveland 1979) estimates $m(x_0) = E[y | x = x_0]$ locally:
\begin{enumerate}
  \item Selects the $f \cdot n$ nearest neighbors of $x_0$
  \item Applies tricubic weights: $w_i = (1 - |d_i|^3)^3$
  \item Fits a weighted regression (linear or constant)
  \item Iterates with robust reweighting (optional)
\end{enumerate}

% ─────────────────────────────────────────────────────────────────────────────
\section{DGP Verification}

DGP: heteroscedastic error $\varepsilon_i \propto (1 + x_i)$; small
coefficient on $t$ (0{,}01). $N = 300$ observations.

\begin{lstlisting}[caption={WLS, Rolling OLS, LOWESS}]
seed(42)
generate df w = exp(-0.002 * t)        // weight: more recent observations more relevant
generate df e = (rnormal()-0.5)*2*(1+x)  // heteroscedasticity in x
generate df y = 1.0 + 2.0*x + 0.01*t + e

let m_ols = ols(y ~ x, df)
let m_wls = wls(y ~ x, df, weights="w")
let m_rol = rols(y ~ x + t, df, window=50)
lowess(df, y, x, frac=0.3)
\end{lstlisting}

\begin{lstlisting}[language={}, basicstyle=\ttfamily\small, frame=none,
  numbers=none, backgroundcolor=\color{codbg}]
OLS:  const=2.362, x=2.263***  (ignores heteroscedasticity)
WLS:  const=2.212, x=2.257***  (weights correct growing variance)
Rolling OLS (last window t=251..300):
  const=1.898, x=2.411, t=0.006
LOWESS: frac=0.30, 300 obs  →  y_hat_lowess added to df
\end{lstlisting}

WLS changes the coefficients only slightly here — temporal weights compensate for
the error growth in $x$ only partially. Rolling OLS reveals that the
coefficient of $x$ in the last window (2{,}41) is slightly higher, suggesting
mild non-stationarity.

% ─────────────────────────────────────────────────────────────────────────────
\section{Syntax}

\begin{lstlisting}
wls(y ~ x, df, weights="weight_column")      // WLS
rols(y ~ x + t, df, window=50)               // Rolling OLS
rolling(y ~ x + t, df, window=50)            // alias
lowess(df, y_var, x_var, frac=0.67, it=3)   // LOWESS
\end{lstlisting}

\begin{nota}
  LOWESS receives the dependent variable \textbf{before} the
  independent one: \texttt{lowess(df, y, x)}. This differs from
  conventional regressions. The parameter \texttt{frac} (fraction of neighbors used at
  each point) controls smoothing: smaller values $\to$ noisier curve;
  larger values $\to$ smoother.
\end{nota}
